\chapter{METHODOLOGY}

\section{System Overview}
Mountain peaks are permanent geographic features whose shapes, as seen from a given coordinate, are unique. A photograph taken in the Khumbu region encodes the observer's position in the geometry of the visible skyline. This project reconstructs that geometry and matches it against a pre-computed reference database to recover the camera's location and heading.

The system has three stages. First, a reference database of horizon profiles is synthesized from the Copernicus GLO-30 DEM using ray-casting. Second, a neural network segments the sky boundary from a query photograph. Third, a two-phase matching engine aligns the extracted profile against the database to determine coordinates.

\section{Horizon Database Generation}
We use the HORAYZON ray-casting library \cite{Steger2022} to compute the horizon profile at each viewpoint in our study region. Each profile is a 1D vector of elevation angles:
\begin{equation}
    V = \{h_1, h_2, \dots, h_{1440}\}
\end{equation}
sampled at $0.25^\circ$ azimuth intervals across a full $360^\circ$ sweep. HORAYZON includes an optional correction for atmospheric refraction, which we enable during generation.

Fog and haze are common in the Khumbu valleys and can obscure distant peaks entirely. A single horizon computed at one fixed search radius would fail under these conditions. We therefore compute three horizons per viewpoint at different maximum search radii---global, local, and restricted---so that the matching engine can select the profile that best reflects the observer's actual visibility. The three tiers are stored for each viewpoint and evaluated during matching.

\begin{figure}[htbp]
    \centering
    \resizebox{\textwidth}{!}{\documentclass[11pt, tikz, border=10pt]{standalone}
\usepackage{pgfplots}
\pgfplotsset{compat=1.18}
\usepackage{amsmath}
\usetikzlibrary{calc, arrows.meta, shadows}

\definecolor{bglight}{HTML}{F8FAFC}
\definecolor{borderc}{HTML}{E2E8F0}
\definecolor{textmain}{HTML}{0F172A}
\definecolor{textsub}{HTML}{475569}
\definecolor{valleygreen}{RGB}{35,95,25}
\definecolor{midgreen}{RGB}{95,155,100}
\definecolor{dirtbrown}{RGB}{185,150,110}
\definecolor{rockbrown}{RGB}{215,200,185}
\definecolor{snowwhite}{RGB}{250,250,250}
\definecolor{queryred}{HTML}{EF4444}
\definecolor{chartblue}{HTML}{0284C7}

\begin{document} 

\begin{tikzpicture}[
    line join=round, 
    line cap=round,
    declare function={
        horizonrad1(\t) = 2.13137 - 0.13137*cos(4*(\t - 70));
        horizonrad2(\t) = 1.80000 - 0.10000*cos(4*(\t - 70));
        horizonrad3(\t) = 1.50000 - 0.08000*cos(4*(\t - 70));
        terrain_base(\u,\v) = 
            0.7*(\u^2 + \v^2)*exp(-0.18*(\u^2 + \v^2))
            + 2.0*exp(-0.8*(\u+1.6)^2 - 0.8*(\v-1.6)^2)
            + 2.2*exp(-0.6*(\u-1.6)^2 - 0.6*(\v+1.6)^2)
            + 1.5*exp(-0.8*(\u+1.6)^2 - 0.8*(\v+1.6)^2)
            + 1.8*exp(-0.8*(\u-1.6)^2 - 0.8*(\v-1.6)^2)
            + 0.15*cos(deg(\u*3))*cos(deg(\v*3));
        % Rotate coordinate system of terrain by 20 degrees
        terrain(\x,\y) = terrain_base(0.9397*\x - 0.3420*\y, 0.3420*\x + 0.9397*\y);
        % Mathematical mapping of 1D curve to the height coordinates of the 2D synthetic box
        plot_y(\t) = -3.0 + 0.71875 * (terrain(horizonrad3(\t)*cos(\t), horizonrad3(\t)*sin(\t)) - 1.0);
    }
]

    % ========================================================
    % 1. TOP ROW: GLOBAL HORIZON
    % ========================================================
    \begin{scope}[yshift=14cm]
        % Column Headers
        \node[font=\large\sffamily\bfseries, text=textmain, anchor=south] at (5.45, 8.2) {Digital Elevation Model};
        \node[font=\large\sffamily\bfseries, text=textmain, anchor=south] at (15.2, 8.2) {Expected Skylines};
        \node[font=\large\sffamily\bfseries, text=textmain, anchor=south] at (25.5, 8.2) {Vector Index};
        
        % Rotated Side Label
        \node[font=\large\sffamily\bfseries, text=textmain, align=center, anchor=south, rotate=90] at (-0.8, 4.5) {Global Horizon\\(Clear Visibility)};

        \begin{axis}[
            width=10.5cm, height=9cm,
            at={(0.2cm, 1.0cm)}, anchor=south west,
            view={-45}{45}, hide axis,
            colormap={terrain}{
                color(0cm)=(valleygreen); color(2cm)=(midgreen);
                color(4cm)=(dirtbrown); color(6cm)=(rockbrown);
                color(8cm)=(snowwhite);
            },
            domain=-3.5:3.5, y domain=-3.5:3.5, samples=28, z buffer=sort,
            point meta min=0.3, point meta max=3.6,
        ]
            \addplot3[only marks, mark=*, mark options={fill=queryred, draw=queryred, line width=0pt}, mark size=1.1pt, domain=45:225, y domain=0:0, samples=160, samples y=1] 
                ({horizonrad1(x)*cos(x)}, {horizonrad1(x)*sin(x)}, {terrain(horizonrad1(x)*cos(x), horizonrad1(x)*sin(x)) + 0.03});

            \addplot3[only marks, mark=*, mark options={fill=textsub!70, draw=none}, mark size=0.8pt, domain=0:1, samples=9, samples y=1] 
                ({x * horizonrad1(90)*cos(90)}, {x * horizonrad1(90)*sin(90)}, {1.3 + x * (terrain(horizonrad1(90)*cos(90), horizonrad1(90)*sin(90)) + 0.03 - 1.3)});
            \addplot3[only marks, mark=*, mark options={fill=textsub!70, draw=none}, mark size=0.8pt, domain=0:1, samples=9, samples y=1] 
                ({x * horizonrad1(180)*cos(180)}, {x * horizonrad1(180)*sin(180)}, {1.3 + x * (terrain(horizonrad1(180)*cos(180), horizonrad1(180)*sin(180)) + 0.03 - 1.3)});

            \addplot3[mark=*, mark options={fill=white, draw=queryred!80, line width=1.2pt}, mark size=3pt] coordinates {({horizonrad1(90)*cos(90)}, {horizonrad1(90)*sin(90)}, {terrain(horizonrad1(90)*cos(90), horizonrad1(90)*sin(90)) + 0.03})};
            \addplot3[mark=*, mark options={fill=white, draw=queryred!80, line width=1.2pt}, mark size=3pt] coordinates {({horizonrad1(180)*cos(180)}, {horizonrad1(180)*sin(180)}, {terrain(horizonrad1(180)*cos(180), horizonrad1(180)*sin(180)) + 0.03})};

            \addplot3[surf, faceted color=black!20, line width=0.3pt] { terrain(x, y) };

            \addplot3[only marks, mark=*, mark options={fill=queryred, draw=queryred, line width=0pt}, mark size=1.1pt, domain=225:405, y domain=0:0, samples=160, samples y=1] 
                ({horizonrad1(x)*cos(x)}, {horizonrad1(x)*sin(x)}, {terrain(horizonrad1(x)*cos(x), horizonrad1(x)*sin(x)) + 0.03});

            \addplot3[only marks, mark=*, mark options={fill=textsub!70, draw=none}, mark size=0.8pt, domain=0:1, samples=9, samples y=1] 
                ({x * horizonrad1(0)*cos(0)}, {x * horizonrad1(0)*sin(0)}, {1.3 + x * (terrain(horizonrad1(0)*cos(0), horizonrad1(0)*sin(0)) + 0.03 - 1.3)});
            \addplot3[only marks, mark=*, mark options={fill=textsub!70, draw=none}, mark size=0.8pt, domain=0:1, samples=9, samples y=1] 
                ({x * horizonrad1(270)*cos(270)}, {x * horizonrad1(270)*sin(270)}, {1.3 + x * (terrain(horizonrad1(270)*cos(270), horizonrad1(270)*sin(270)) + 0.03 - 1.3)});

            \addplot3[mark=*, mark options={fill=white, draw=queryred!80, line width=1.2pt}, mark size=3pt] coordinates {({horizonrad1(0)*cos(0)}, {horizonrad1(0)*sin(0)}, {terrain(horizonrad1(0)*cos(0), horizonrad1(0)*sin(0)) + 0.03})};
            \addplot3[mark=*, mark options={fill=white, draw=queryred!80, line width=1.2pt}, mark size=3pt] coordinates {({horizonrad1(270)*cos(270)}, {horizonrad1(270)*sin(270)}, {terrain(horizonrad1(270)*cos(270), horizonrad1(270)*sin(270)) + 0.03})};

            \node[font=\scriptsize\sffamily\bfseries, text=textsub, anchor=north, yshift=-2mm] at (axis cs: {horizonrad1(0)*cos(0)}, {horizonrad1(0)*sin(0)}, {terrain(horizonrad1(0)*cos(0), horizonrad1(0)*sin(0)) + 0.03}) {$0^\circ$};
            \node[font=\scriptsize\sffamily\bfseries, text=textsub, anchor=south, yshift=2.5mm] at (axis cs: {horizonrad1(90)*cos(90)}, {horizonrad1(90)*sin(90)}, {terrain(horizonrad1(90)*cos(90), horizonrad1(90)*sin(90)) + 0.03}) {$90^\circ$};
            \node[font=\scriptsize\sffamily\bfseries, text=textsub, anchor=south, yshift=2.5mm] at (axis cs: {horizonrad1(180)*cos(180)}, {horizonrad1(180)*sin(180)}, {terrain(horizonrad1(180)*cos(180), horizonrad1(180)*sin(180)) + 0.03}) {$180^\circ$};
            \node[font=\scriptsize\sffamily\bfseries, text=textsub, anchor=north, yshift=-2mm] at (axis cs: {horizonrad1(270)*cos(270)}, {horizonrad1(270)*sin(270)}, {terrain(horizonrad1(270)*cos(270), horizonrad1(270)*sin(270)) + 0.03}) {$270^\circ$};

            \coordinate (ObserverPos1) at (axis cs:0,0,1.3);
        \end{axis}

        % Restored Viewpoint Pin
        \node[anchor=south, inner sep=0, yshift=-0.4mm] at (ObserverPos1) {
            \begin{tikzpicture}[scale=0.14]
                \fill[black!15] (0.08,-0.08) circle (0.3);
                \fill[chartblue, draw=white, line width=1.1pt] (0,0) .. controls (-1.2,1.8) and (-1.8,3.6) .. (0,5.4) .. controls (1.8,3.6) and (1.2,1.8) .. (0,0) -- cycle;
                \fill[white] (0,3.6) circle (0.6);
            \end{tikzpicture}
        };

        \begin{scope}[shift={(11.2,2.2)}]
            \begin{axis}[
                width=8cm, height=6.5cm, anchor=south west, axis lines=left, axis line style={textmain, line width=1.5pt},
                xmin=0, xmax=360, ymin=1.0, ymax=4.2, xtick={0, 90, 180, 270, 360},
                xticklabels={0$^\circ$, 90$^\circ$, 180$^\circ$, 270$^\circ$, 360$^\circ$}, xticklabel style={font=\scriptsize\sffamily, text=textsub}, ytick=\empty,
                xlabel={Azimuth Angle ($\theta$)}, ylabel={Maximum Elevation seen from Viewpoint},
                xlabel style={font=\small\sffamily, text=textmain, yshift=-1mm}, ylabel style={font=\small\sffamily, text=textmain, yshift=2mm},
            ]
                \addplot[queryred, line width=2.5pt, domain=0:360, samples=150, smooth] {terrain(horizonrad1(x)*cos(x), horizonrad1(x)*sin(x))};
            \end{axis}
        \end{scope}
        
        % Vector Symbol [ ... ]
        \node (vec1) at (20.5, 5.0) [xscale=1.5, yscale=3.5, font=\large] {$\begin{bmatrix} \textcolor{queryred}{\vdots} \end{bmatrix}$};
    \end{scope}

    % ========================================================
    % 2. MIDDLE ROW: LOCAL HORIZON
    % ========================================================
    \begin{scope}[yshift=7cm]
        \node[font=\large\sffamily\bfseries, text=textmain, align=center, anchor=south, rotate=90] at (-0.8, 4.5) {Local Horizon\\(Moderate Haze)};

        \begin{axis}[
            width=10.5cm, height=9cm,
            at={(0.2cm, 1.0cm)}, anchor=south west,
            view={-45}{45}, hide axis,
            colormap={terrain}{
                color(0cm)=(valleygreen); color(2cm)=(midgreen);
                color(4cm)=(dirtbrown); color(6cm)=(rockbrown);
                color(8cm)=(snowwhite);
            },
            domain=-3.5:3.5, y domain=-3.5:3.5, samples=28, z buffer=sort,
            point meta min=0.3, point meta max=3.6,
        ]
            \addplot3[only marks, mark=*, mark options={fill=queryred, draw=queryred, line width=0pt}, mark size=1.1pt, domain=45:225, y domain=0:0, samples=160, samples y=1] 
                ({horizonrad2(x)*cos(x)}, {horizonrad2(x)*sin(x)}, {terrain(horizonrad2(x)*cos(x), horizonrad2(x)*sin(x)) + 0.03});

            \addplot3[only marks, mark=*, mark options={fill=textsub!70, draw=none}, mark size=0.8pt, domain=0:1, samples=9, samples y=1] 
                ({x * horizonrad2(90)*cos(90)}, {x * horizonrad2(90)*sin(90)}, {1.3 + x * (terrain(horizonrad2(90)*cos(90), horizonrad2(90)*sin(90)) + 0.03 - 1.3)});
            \addplot3[only marks, mark=*, mark options={fill=textsub!70, draw=none}, mark size=0.8pt, domain=0:1, samples=9, samples y=1] 
                ({x * horizonrad2(180)*cos(180)}, {x * horizonrad2(180)*sin(180)}, {1.3 + x * (terrain(horizonrad2(180)*cos(180), horizonrad2(180)*sin(180)) + 0.03 - 1.3)});

            \addplot3[mark=*, mark options={fill=white, draw=queryred!80, line width=1.2pt}, mark size=3pt] coordinates {({horizonrad2(90)*cos(90)}, {horizonrad2(90)*sin(90)}, {terrain(horizonrad2(90)*cos(90), horizonrad2(90)*sin(90)) + 0.03})};
            \addplot3[mark=*, mark options={fill=white, draw=queryred!80, line width=1.2pt}, mark size=3pt] coordinates {({horizonrad2(180)*cos(180)}, {horizonrad2(180)*sin(180)}, {terrain(horizonrad2(180)*cos(180), horizonrad2(180)*sin(180)) + 0.03})};

            \addplot3[surf, faceted color=black!20, line width=0.3pt] { terrain(x, y) };

            \addplot3[only marks, mark=*, mark options={fill=queryred, draw=queryred, line width=0pt}, mark size=1.1pt, domain=225:405, y domain=0:0, samples=160, samples y=1] 
                ({horizonrad2(x)*cos(x)}, {horizonrad2(x)*sin(x)}, {terrain(horizonrad2(x)*cos(x), horizonrad2(x)*sin(x)) + 0.03});

            \addplot3[only marks, mark=*, mark options={fill=textsub!70, draw=none}, mark size=0.8pt, domain=0:1, samples=9, samples y=1] 
                ({x * horizonrad2(0)*cos(0)}, {x * horizonrad2(0)*sin(0)}, {1.3 + x * (terrain(horizonrad2(0)*cos(0), horizonrad2(0)*sin(0)) + 0.03 - 1.3)});
            \addplot3[only marks, mark=*, mark options={fill=textsub!70, draw=none}, mark size=0.8pt, domain=0:1, samples=9, samples y=1] 
                ({x * horizonrad2(270)*cos(270)}, {x * horizonrad2(270)*sin(270)}, {1.3 + x * (terrain(horizonrad2(270)*cos(270), horizonrad2(270)*sin(270)) + 0.03 - 1.3)});

            \addplot3[mark=*, mark options={fill=white, draw=queryred!80, line width=1.2pt}, mark size=3pt] coordinates {({horizonrad2(0)*cos(0)}, {horizonrad2(0)*sin(0)}, {terrain(horizonrad2(0)*cos(0), horizonrad2(0)*sin(0)) + 0.03})};
            \addplot3[mark=*, mark options={fill=white, draw=queryred!80, line width=1.2pt}, mark size=3pt] coordinates {({horizonrad2(270)*cos(270)}, {horizonrad2(270)*sin(270)}, {terrain(horizonrad2(270)*cos(270), horizonrad2(270)*sin(270)) + 0.03})};

            \node[font=\scriptsize\sffamily\bfseries, text=textsub, anchor=north, yshift=-2mm] at (axis cs: {horizonrad2(0)*cos(0)}, {horizonrad2(0)*sin(0)}, {terrain(horizonrad2(0)*cos(0), horizonrad2(0)*sin(0)) + 0.03}) {$0^\circ$};
            \node[font=\scriptsize\sffamily\bfseries, text=textsub, anchor=south, yshift=2.5mm] at (axis cs: {horizonrad2(90)*cos(90)}, {horizonrad2(90)*sin(90)}, {terrain(horizonrad2(90)*cos(90), horizonrad2(90)*sin(90)) + 0.03}) {$90^\circ$};
            \node[font=\scriptsize\sffamily\bfseries, text=textsub, anchor=south, yshift=2.5mm] at (axis cs: {horizonrad2(180)*cos(180)}, {horizonrad2(180)*sin(180)}, {terrain(horizonrad2(180)*cos(180), horizonrad2(180)*sin(180)) + 0.03}) {$180^\circ$};
            \node[font=\scriptsize\sffamily\bfseries, text=textsub, anchor=north, yshift=-2mm] at (axis cs: {horizonrad2(270)*cos(270)}, {horizonrad2(270)*sin(270)}, {terrain(horizonrad2(270)*cos(270), horizonrad2(270)*sin(270)) + 0.03}) {$270^\circ$};

            \coordinate (ObserverPos2) at (axis cs:0,0,1.3);
        \end{axis}

        % Restored Viewpoint Pin
        \node[anchor=south, inner sep=0, yshift=-0.4mm] at (ObserverPos2) {
            \begin{tikzpicture}[scale=0.14]
                \fill[black!15] (0.08,-0.08) circle (0.3);
                \fill[chartblue, draw=white, line width=1.1pt] (0,0) .. controls (-1.2,1.8) and (-1.8,3.6) .. (0,5.4) .. controls (1.8,3.6) and (1.2,1.8) .. (0,0) -- cycle;
                \fill[white] (0,3.6) circle (0.6);
            \end{tikzpicture}
        };

        \begin{scope}[shift={(11.2,2.2)}]
            \begin{axis}[
                width=8cm, height=6.5cm, anchor=south west, axis lines=left, axis line style={textmain, line width=1.5pt},
                xmin=0, xmax=360, ymin=1.0, ymax=4.2, xtick={0, 90, 180, 270, 360},
                xticklabels={0$^\circ$, 90$^\circ$, 180$^\circ$, 270$^\circ$, 360$^\circ$}, xticklabel style={font=\scriptsize\sffamily, text=textsub}, ytick=\empty,
                xlabel={Azimuth Angle ($\theta$)}, ylabel={Maximum Elevation seen from Viewpoint},
                xlabel style={font=\small\sffamily, text=textmain, yshift=-1mm}, ylabel style={font=\small\sffamily, text=textmain, yshift=2mm},
            ]
                \addplot[queryred, line width=2.5pt, domain=0:360, samples=150, smooth] {terrain(horizonrad2(x)*cos(x), horizonrad2(x)*sin(x))};
            \end{axis}
        \end{scope}
        
        % Vector Symbol [ ... ]
        \node (vec2) at (20.5, 5.0) [xscale=1.5, yscale=3.5, font=\large] {$\begin{bmatrix} \textcolor{queryred}{\vdots} \end{bmatrix}$};
    \end{scope}

    % ========================================================
    % 3. BOTTOM ROW: RESTRICTED HORIZON
    % ========================================================
    \begin{scope}[yshift=0cm]
        \node[font=\large\sffamily\bfseries, text=textmain, align=center, anchor=south, rotate=90] at (-0.8, 4.5) {Restricted Horizon\\(Dense Fog)};

        \begin{axis}[
            width=10.5cm, height=9cm,
            at={(0.2cm, 1.0cm)}, anchor=south west,
            view={-45}{45}, hide axis,
            colormap={terrain}{
                color(0cm)=(valleygreen); color(2cm)=(midgreen);
                color(4cm)=(dirtbrown); color(6cm)=(rockbrown);
                color(8cm)=(snowwhite);
            },
            domain=-3.5:3.5, y domain=-3.5:3.5, samples=28, z buffer=sort,
            point meta min=0.3, point meta max=3.6,
        ]
            \addplot3[only marks, mark=*, mark options={fill=queryred, draw=queryred, line width=0pt}, mark size=1.1pt, domain=45:225, y domain=0:0, samples=160, samples y=1] 
                ({horizonrad3(x)*cos(x)}, {horizonrad3(x)*sin(x)}, {terrain(horizonrad3(x)*cos(x), horizonrad3(x)*sin(x)) + 0.03});

            \addplot3[only marks, mark=*, mark options={fill=textsub!70, draw=none}, mark size=0.8pt, domain=0:1, samples=9, samples y=1] 
                ({x * horizonrad3(90)*cos(90)}, {x * horizonrad3(90)*sin(90)}, {1.3 + x * (terrain(horizonrad3(90)*cos(90), horizonrad3(90)*sin(90)) + 0.03 - 1.3)});
            \addplot3[only marks, mark=*, mark options={fill=textsub!70, draw=none}, mark size=0.8pt, domain=0:1, samples=9, samples y=1] 
                ({x * horizonrad3(180)*cos(180)}, {x * horizonrad3(180)*sin(180)}, {1.3 + x * (terrain(horizonrad3(180)*cos(180), horizonrad3(180)*sin(180)) + 0.03 - 1.3)});

            \addplot3[mark=*, mark options={fill=white, draw=queryred!80, line width=1.2pt}, mark size=3pt] coordinates {({horizonrad3(90)*cos(90)}, {horizonrad3(90)*sin(90)}, {terrain(horizonrad3(90)*cos(90), horizonrad3(90)*sin(90)) + 0.03})};
            \addplot3[mark=*, mark options={fill=white, draw=queryred!80, line width=1.2pt}, mark size=3pt] coordinates {({horizonrad3(180)*cos(180)}, {horizonrad3(180)*sin(180)}, {terrain(horizonrad3(180)*cos(180), horizonrad3(180)*sin(180)) + 0.03})};

            \addplot3[surf, faceted color=black!20, line width=0.3pt] { terrain(x, y) };

            \addplot3[only marks, mark=*, mark options={fill=queryred, draw=queryred, line width=0pt}, mark size=1.1pt, domain=225:405, y domain=0:0, samples=160, samples y=1] 
                ({horizonrad3(x)*cos(x)}, {horizonrad3(x)*sin(x)}, {terrain(horizonrad3(x)*cos(x), horizonrad3(x)*sin(x)) + 0.03});

            \addplot3[only marks, mark=*, mark options={fill=textsub!70, draw=none}, mark size=0.8pt, domain=0:1, samples=9, samples y=1] 
                ({x * horizonrad3(0)*cos(0)}, {x * horizonrad3(0)*sin(0)}, {1.3 + x * (terrain(horizonrad3(0)*cos(0), horizonrad3(0)*sin(0)) + 0.03 - 1.3)});
            \addplot3[only marks, mark=*, mark options={fill=textsub!70, draw=none}, mark size=0.8pt, domain=0:1, samples=9, samples y=1] 
                ({x * horizonrad3(270)*cos(270)}, {x * horizonrad3(270)*sin(270)}, {1.3 + x * (terrain(horizonrad3(270)*cos(270), horizonrad3(270)*sin(270)) + 0.03 - 1.3)});

            \addplot3[mark=*, mark options={fill=white, draw=queryred!80, line width=1.2pt}, mark size=3pt] coordinates {({horizonrad3(0)*cos(0)}, {horizonrad3(0)*sin(0)}, {terrain(horizonrad3(0)*cos(0), horizonrad3(0)*sin(0)) + 0.03})};
            \addplot3[mark=*, mark options={fill=white, draw=queryred!80, line width=1.2pt}, mark size=3pt] coordinates {({horizonrad3(270)*cos(270)}, {horizonrad3(270)*sin(270)}, {terrain(horizonrad3(270)*cos(270), horizonrad3(270)*sin(270)) + 0.03})};

            \node[font=\scriptsize\sffamily\bfseries, text=textsub, anchor=north, yshift=-2mm] at (axis cs: {horizonrad3(0)*cos(0)}, {horizonrad3(0)*sin(0)}, {terrain(horizonrad3(0)*cos(0), horizonrad3(0)*sin(0)) + 0.03}) {$0^\circ$};
            \node[font=\scriptsize\sffamily\bfseries, text=textsub, anchor=south, yshift=2.5mm] at (axis cs: {horizonrad3(90)*cos(90)}, {horizonrad3(90)*sin(90)}, {terrain(horizonrad3(90)*cos(90), horizonrad3(90)*sin(90)) + 0.03}) {$90^\circ$};
            \node[font=\scriptsize\sffamily\bfseries, text=textsub, anchor=south, yshift=2.5mm] at (axis cs: {horizonrad3(180)*cos(180)}, {horizonrad3(180)*sin(180)}, {terrain(horizonrad3(180)*cos(180), horizonrad3(180)*sin(180)) + 0.03}) {$180^\circ$};
            \node[font=\scriptsize\sffamily\bfseries, text=textsub, anchor=north, yshift=-2mm] at (axis cs: {horizonrad3(270)*cos(270)}, {horizonrad3(270)*sin(270)}, {terrain(horizonrad3(270)*cos(270), horizonrad3(270)*sin(270)) + 0.03}) {$270^\circ$};

            \coordinate (ObserverPos3) at (axis cs:0,0,1.3);
        \end{axis}

        % Restored Viewpoint Pin
        \node[anchor=south, inner sep=0, yshift=-0.4mm] at (ObserverPos3) {
            \begin{tikzpicture}[scale=0.14]
                \fill[black!15] (0.08,-0.08) circle (0.3);
                \fill[chartblue, draw=white, line width=1.1pt] (0,0) .. controls (-1.2,1.8) and (-1.8,3.6) .. (0,5.4) .. controls (1.8,3.6) and (1.2,1.8) .. (0,0) -- cycle;
                \fill[white] (0,3.6) circle (0.6);
            \end{tikzpicture}
        };

        \begin{scope}[shift={(11.2,2.2)}]
            \begin{axis}[
                width=8cm, height=6.5cm, anchor=south west, axis lines=left, axis line style={textmain, line width=1.5pt},
                xmin=0, xmax=360, ymin=1.0, ymax=4.2, xtick={0, 90, 180, 270, 360},
                xticklabels={0$^\circ$, 90$^\circ$, 180$^\circ$, 270$^\circ$, 360$^\circ$}, xticklabel style={font=\scriptsize\sffamily, text=textsub}, ytick=\empty,
                xlabel={Azimuth Angle ($\theta$)}, ylabel={Maximum Elevation seen from Viewpoint},
                xlabel style={font=\small\sffamily, text=textmain, yshift=-1mm}, ylabel style={font=\small\sffamily, text=textmain, yshift=2mm},
            ]
                \addplot[queryred, line width=2.5pt, domain=0:360, samples=150, smooth] {terrain(horizonrad3(x)*cos(x), horizonrad3(x)*sin(x))};
            \end{axis}
        \end{scope}
        
        % Vector Symbol [ ... ]
        \node (vec3) at (20.5, 5.0) [xscale=1.5, yscale=3.5, font=\large] {$\begin{bmatrix} \textcolor{queryred}{\vdots} \end{bmatrix}$};

        % ========================================================
        % GENERATIVE SYNTHETIC DATA BLOCK (Below 3rd Row 1D Graph)
        % ========================================================
        
        % Sleek slate-gray arrow moved cleanly below the xlabel text to avoid clutter/overlap
        \draw[textsub, line width=1.8pt, -{Stealth[length=3mm, width=2mm]}] (15.2, 0.0) -- (15.2, -0.6);
        
        % Box 3 (deepest background card - only outline visible on bottom/right)
        \draw[textsub, thick, rounded corners=4pt, fill=white] (11.5, -3.3) rectangle (19.5, -1.0);
        
        % Box 2 (middle background card - only outline visible on bottom/right)
        \draw[textsub, thick, rounded corners=4pt, fill=white] (11.35, -3.15) rectangle (19.35, -0.85);

        % Hidden PGFPlots axis precisely matching 1D graph coordinate space
        \begin{axis}[
            width=8cm, height=2.3cm,
            at={(11.2cm, -3.0cm)}, anchor=south west,
            xmin=0, xmax=360, ymin=1.0, ymax=4.2,
            scale only axis, % Forces coordinate viewport to fill the specified box boundaries
            hide axis
        ]
            % Perfect clipping container for the rounded corners
            \begin{scope}
                \clip[rounded corners=4pt] (0, 1.0) rectangle (360, 4.2);
                
                % High-altitude Himalayan sky background (now completely fills the box area)
                \shade[top color=chartblue!90!black, bottom color=chartblue!15!white] (0, 1.0) rectangle (360, 4.2);
                
                % Snow-capped mountain silhouette matching the red 1D curve 1-to-1 (now completely fills the box area)
                \shade[top color=snowwhite, bottom color=textmain!85!black, draw=none] 
                    (0, 1.0) -- plot[domain=0:360, samples=80, smooth] (\x, {terrain(horizonrad3(\x)*cos(\x), horizonrad3(\x)*sin(\x))}) -- (360, 1.0) -- cycle;
            \end{scope}
        \end{axis}
        
        % Protective border matching the background box (Main Box on top)
        \draw[textsub, thick, rounded corners=4pt] (11.2, -3.0) rectangle (19.2, -0.7);
        
        % Synthetic Data label below the generated stacked frame
        \node[font=\small\sffamily\bfseries, text=textmain, anchor=north] at (15.35, -3.4) {Simulated Ground-Level Imagery};
        
    \end{scope}

    % ========================================================
    % 4. SHARED CENTRAL GRID & CONNECTIONS
    % ========================================================
    \begin{scope}[shift={(23.2, 8.5)}]
        % Manual Grid (Outer boundary + clean internal divisions)
        \draw[borderc, thick, fill=white] (0,0) rectangle (4.0,4.0);
        \draw[borderc, thin, step=0.8cm] (0,0) grid (4.0,4.0);
        
        % Highlighted cell (Center area)
        \fill[queryred!20] (1.6,1.6) rectangle (2.4,2.4);

        \coordinate (cell0) at (1.8, 1.8); % The center cell itself
        \coordinate (cell1) at (2.0, 2.0); % The center cell itself
        \coordinate (cell2) at (1.8, 2.2); % The center cell itself
        \draw[textmain, ultra thick, -{Stealth[length=3mm]}] (vec1.east) to[out=0, in=135] (cell2);
        \draw[textmain, ultra thick, -{Stealth[length=3mm]}] (vec2.east) to[out=0, in=180] (cell1);
        \draw[textmain, ultra thick, -{Stealth[length=3mm]}] (vec3.east) to[out=0, in=-135] (cell0);

        
        % Teardrop Pin added on top of the highlighted grid cell
        \node[anchor=south, inner sep=0, yshift=0.5mm] at (2.0, 2.0) {
            \begin{tikzpicture}[scale=0.14]
                \fill[black!15] (0.08,-0.08) circle (0.3);
                \fill[chartblue, draw=white, line width=1.1pt] (0,0) .. controls (-1.2,1.8) and (-1.8,3.6) .. (0,5.4) .. controls (1.8,3.6) and (1.2,1.8) .. (0,0) -- cycle;
                \fill[white] (0,3.6) circle (0.6);
            \end{tikzpicture}
        };
        
    \end{scope}

\end{tikzpicture}
\end{document}
}
    \caption{Raycasting to find the expected horizons at a given viewpoint using progressive search radii.}
    \label{fig:design}
\end{figure}

\section{Skyline Extraction}
We use a U-Net with a MobileNet-V3 encoder to segment query photographs into sky and terrain \cite{long2015fully, howard2019mobilenetv3}. The model is pre-trained on GeoPose3K \cite{Brejcha2017}, which provides over 3,000 labeled alpine mountain images, and then fine-tuned on synthetic renderings of the Khumbu region generated from the DEM. These synthetic images are augmented with procedural fog textures to better represent local weather. A Canny edge detector \cite{Canny1986} serves as a fallback when the neural network produces low-confidence output.

\begin{figure}[htbp]
    \centering
    \resizebox{\textwidth}{!}{\documentclass[11pt, tikz, border=10pt]{standalone}
\usepackage{pgfplots}
\pgfplotsset{compat=1.18}
\usetikzlibrary{calc, arrows.meta, positioning, shapes.geometric, fit}

\definecolor{bglight}{HTML}{F8FAFC}
\definecolor{borderc}{HTML}{E2E8F0}
\definecolor{textmain}{HTML}{0F172A}
\definecolor{textsub}{HTML}{475569}
\definecolor{chartblue}{HTML}{0284C7}
\definecolor{queryred}{HTML}{EF4444}
\definecolor{valleygreen}{RGB}{35,95,25}
\definecolor{midgreen}{RGB}{95,155,100}
\definecolor{dirtbrown}{RGB}{185,150,110}
\definecolor{rockbrown}{RGB}{215,200,185}
\definecolor{snowwhite}{RGB}{250,250,250}

\begin{document}

\begin{tikzpicture}[
    line join=round, 
    line cap=round,
    % Node styles
    boxstyle/.style={
        draw=borderc, thick, fill=white, rounded corners=4pt, 
        align=center, text=textmain, font=\sffamily\small,
        minimum width=2.6cm, minimum height=1.3cm,
        inner sep=6pt
    },
    datastyle/.style={
        draw=borderc, thick, fill=bglight, rounded corners=3pt,
        align=center, text=textmain, font=\sffamily\footnotesize,
        minimum width=2.4cm, minimum height=1.2cm,
        inner sep=6pt
    },
    arrowstyle/.style={
        -{Stealth[length=3.5mm, width=2.2mm]},
        line width=2pt,
        draw=textsub
    },
    skipstyle/.style={
        dashed,
        -{Stealth[length=2.5mm]},
        line width=1.2pt,
        draw=chartblue!80!black
    },
    fallbackstyle/.style={
        dashed,
        -{Stealth[length=3.5mm, width=2.2mm]},
        line width=1.8pt,
        draw=textsub!60
    },
    % Mathematical function representing the exact geographical skyline for the user image pipeline
    declare function={
        user_skyline(\x) = 0.35 + 0.20*sin(\x*200 + 45) + 0.08*cos(\x*380);
    }
]

    % ========================================================
    % SECTION A: TRAINING DATASET PIPELINE (Left Side)
    % ========================================================
    \node[font=\large\sffamily\bfseries, text=textmain] (title1) at (1.8, 9.5) {1. Training Dataset Generation};

    % Real Data Pipeline (Photo + Mask)
    \begin{scope}[local bounding box=realdata]
        % Real Photo
        \node[datastyle] (geophoto) at (0, 7.5) {GeoPose3K\\Real Photo};
        \begin{scope}[shift={($(geophoto.south) + (-1.2,-1.3)$)}]
            \draw[borderc, thick, fill=white] (0,0) rectangle (2.4, 1.0);
            \shade[top color=chartblue!30, bottom color=white] (0,0) rectangle (2.4, 1.0);
            \fill[midgreen] (0,0) -- (0, 0.4) to[out=20, in=160] (1.2, 0.3) to[out=-20, in=180] (2.4, 0.1) -- (2.4,0) -- cycle;
            \fill[valleygreen] (0,0) -- (0, 0.2) to[out=10, in=170] (1.6, 0.1) to[out=-10, in=180] (2.4, 0.05) -- (2.4,0) -- cycle;
        \end{scope}

        % Ground Truth Binary Mask
        \node[datastyle] (geomask) at (3.8, 7.5) {Ground Truth\\Sky-Terrain Mask};
        \begin{scope}[shift={($(geomask.south) + (-1.2,-1.3)$)}]
            \draw[borderc, thick, fill=textmain] (0,0) rectangle (2.4, 1.0);
            \fill[white] (0, 1.0) -- (0, 0.4) to[out=20, in=160] (1.2, 0.3) to[out=-20, in=180] (2.4, 0.1) -- (2.4,1.0) -- cycle;
        \end{scope}
    \end{scope}
    
    % Synthetic Data Pipeline (Horizon -> Rendering)
    \begin{scope}[local bounding box=synthdata]
        % Synthetic Renderings (Input to U-Net)
        \node[datastyle] (synthrender_img) at (0, 3.5) {Synthetic\\Rendering};
        \begin{scope}[shift={($(synthrender_img.south) + (-1.2,-1.3)$)}]
            \draw[borderc, thick, fill=white] (0,0) rectangle (2.4, 1.0);
            \shade[top color=chartblue!80!black, bottom color=chartblue!15!white] (0,0) rectangle (2.4, 1.0);
            \shade[top color=snowwhite, bottom color=textmain!80!black, draw=none] 
                (0,0) -- plot[domain=0:2.4, samples=40] (\x, {0.3 + 0.20*sin(\x*360 + 40)}) -- (2.4,0) -- cycle;
        \end{scope}

        % Synthetic Training Masks (Ground Truth for U-Net)
        \node[datastyle] (synthrender_mask) at (3.8, 3.5) {Synthetic\\Binary Mask};
        \begin{scope}[shift={($(synthrender_mask.south) + (-1.2,-1.3)$)}]
            \draw[borderc, thick, fill=textmain] (0,0) rectangle (2.4, 1.0);
            \fill[white] (0, 1.0) -- plot[domain=0:2.4, samples=40] (\x, {0.3 + 0.20*sin(\x*360 + 40)}) -- (2.4,1.0) -- cycle;
        \end{scope}

        \draw[arrowstyle, thin, <->] (synthrender_img) -- (synthrender_mask);
    \end{scope}

    % Group boxes
    \node[draw=textsub!30, dashed, inner sep=12pt, rounded corners=6pt, fit=(realdata)] (boxreal) {};
    \node[draw=textsub!30, dashed, inner sep=12pt, rounded corners=6pt, fit=(synthdata)] (boxsynth) {};
    \node[anchor=north, text=textsub, font=\sffamily\bfseries\scriptsize] at (boxreal.north) {REAL-WORLD DOMAIN};
    \node[anchor=north, text=textsub, font=\sffamily\bfseries\scriptsize] at (boxsynth.north) {SYNTHETIC DOMAIN};

    % ========================================================
    % SECTION B: STANDARD U-NET SCHEMATIC (Center)
    % ========================================================
    \node[font=\large\sffamily\bfseries, text=textmain] (title2) at (11.5, 9.5) {2. U-Net Sky Segmentation};

    \begin{scope}[shift={(8.2, 1.5)}]
        % Contracting Path (Encoder Blocks - kept at uniform 2.6cm)
        \node[boxstyle, fill=chartblue!8, minimum width=2.6cm] (enc1) at (0, 6.0) {Encoder Block 1\\(MobileNet)};
        \node[boxstyle, fill=chartblue!12, minimum width=2.6cm] (enc2) at (0.8, 4.2) {Encoder Block 2\\(MobileNet)};
        \node[boxstyle, fill=chartblue!16, minimum width=2.6cm] (enc3) at (1.6, 2.4) {Encoder Block 3\\(MobileNet)};
        
        % Bottleneck (Shifted horizontally to 3.4 for spacing)
        \node[boxstyle, fill=chartblue!25, minimum width=2.6cm] (bottleneck) at (3.4, 0.6) {Bottleneck\\(Latent Feature)};

        % Expanding Path (Decoder Blocks - shifted horizontally to start at 5.2 for a 1.0cm gap)
        \node[boxstyle, fill=queryred!8, minimum width=2.6cm] (dec3) at (5.2, 2.4) {Decoder Block 3\\(Upsampling)};
        \node[boxstyle, fill=queryred!12, minimum width=2.6cm] (dec2) at (6.0, 4.2) {Decoder Block 2\\(Upsampling)};
        \node[boxstyle, fill=queryred!16, minimum width=2.6cm] (dec1) at (6.8, 6.0) {Decoder Block 1\\(Upsampling)};

        % Encoder Downsampling Arrows (Clean, separated diagonal lines connecting blocks)
        \draw[arrowstyle, thin] (enc1.south) -- (enc2.north);
        \draw[arrowstyle, thin] (enc2.south) -- (enc3.north);
        \draw[arrowstyle, thin] (enc3.south) to[out=210, in=135] (bottleneck.west);

        % Decoder Upsampling Arrows (Clean, separated diagonal lines connecting blocks)
        \draw[arrowstyle, thin] (bottleneck.east) to[out=0, in=-45] (dec3.south);
        \draw[arrowstyle, thin] (dec3.north) -- (dec2.south);
        \draw[arrowstyle, thin] (dec2.north) -- (dec1.south);

        % Skip Connections (All lines, especially enc3 to dec3, are now clearly visible)
        \draw[skipstyle] (enc1.east) -- (dec1.west);
        \draw[skipstyle] (enc2.east) -- (dec2.west);
        \draw[skipstyle] (enc3.east) -- (dec3.west);
    \end{scope}

    % Dataset connection arrows to U-Net training phase (Kept unchanged)
    \draw[arrowstyle, chartblue] (boxreal.east) to[out=70, in=180] (enc1.west);
    \draw[arrowstyle, chartblue] (boxsynth.east) to[out=0, in=180] (enc1.west);

    % ========================================================
    % SECTION C: PIPELINE INFERENCE & EXTRACTION (Right Side)
    % ========================================================
    \node[font=\large\sffamily\bfseries, text=textmain] (title3) at (20.5, 9.5) {3. Skyline Extraction};

    % Inference Nodes - Primary Pipeline
    \node[boxstyle] (userimg) at (19.5, 7.5) {Input\\User Image};
    \node[boxstyle, fill=queryred!10] (segmask) at (19.5, 4.5) {Segmented\\Binary Mask};
    \node[boxstyle, fill=valleygreen!10] (extraction) at (19.5, 1.5) {Extracted\\Skyline};

    % Illustrative Thumbnail Assets (Now sharing a unified mathematical skyline function)
    % User Image Thumbnail
    \begin{scope}[shift={($(userimg.east) + (0.4,-0.5)$)}]
        \draw[borderc, thick, fill=white] (0,0) rectangle (1.8, 1.0);
        \shade[top color=chartblue!40, bottom color=white] (0,0) rectangle (1.8, 1.0);
        \fill[midgreen] (0,0) -- plot[domain=0:1.8, samples=40] (\x, {user_skyline(\x)}) -- (1.8,0) -- cycle;
    \end{scope}

    % Binary Mask Thumbnail
    \begin{scope}[shift={($(segmask.east) + (0.4,-0.5)$)}]
        \draw[borderc, thick, fill=textmain] (0,0) rectangle (1.8, 1.0);
        \fill[white] (0, 1.0) -- plot[domain=0:1.8, samples=40] (\x, {user_skyline(\x)}) -- (1.8,1.0) -- cycle;
    \end{scope}

    % Extracted Skyline Plot Thumbnail
    \begin{scope}[shift={($(extraction.east) + (0.4,-0.5)$)}]
        \draw[borderc, thick, fill=white] (0,0) rectangle (1.8, 1.0);
        \draw[queryred, line width=1.5pt] plot[domain=0:1.8, samples=40] (\x, {user_skyline(\x)});
    \end{scope}

    % Tightened U-Net input arrow (significantly reduced radius to prevent towering over titles)
    \draw[arrowstyle] (userimg.north west) to[out=165, in=15] (enc1.north);
    
    % U-Net Decoder -> Segmented Mask
    \draw[arrowstyle, chartblue] (dec1.east) to[out=0, in=180] (segmask.west);
    
    % Mask -> Extracted Skyline
    \draw[arrowstyle] (segmask) -- (extraction);

    % Direct Canny Edge Detection Fallback Bypass (Rerouted completely to the left of the inference stack)
    \draw[fallbackstyle] (userimg.west) to[out=210, in=190] 
        node[
            font=\sffamily\bfseries\scriptsize,
            text=textsub,
            fill=white,
            inner sep=3.5pt,
            rotate=90,
            anchor=center
        ] {\shortstack{Canny Edge Detection\\(Fallback)}}
    (extraction.west);

\end{tikzpicture}
\end{document}
}
    \caption{Pipeline for skyline extraction from photographs.}
    \label{fig:skyline_extract}
\end{figure}

\section{Matching Engine}
Because the camera heading is unknown, the extracted skyline must be compared against every angular offset of every reference profile in the database. A brute-force search is prohibitively slow, so we split the matching into two phases.

\subsection{Phase 1: Sliding-Window Euclidean Search}
The query profile is treated as a subsequence that slides across each 1440-point reference vector. We compute the Euclidean distance at every offset using FFT-based cross-correlation, which reduces the per-profile cost from $O(Nm)$ to $O(N \log N)$ \cite{mueen2024mass, Rakthanmanon2013}. We deliberately skip z-normalization so that the absolute elevation values---which carry strong discriminative signal in mountain terrain \cite{tonge2014parsing}---are preserved. This stage produces a ranked short-list of candidate viewpoints.

\subsection{Phase 2: DTW Refinement}
Differences in camera focal length and tilt introduce non-linear distortions that Euclidean distance cannot absorb. Dynamic Time Warping handles these elastic deformations by allowing the alignment path to stretch or compress locally. We constrain the computation with Sakoe-Chiba bands and apply early abandoning and pruning to discard poor candidates before the full cost matrix is evaluated \cite{Rakthanmanon2013}.

\begin{figure}[htbp]
    \centering
    \resizebox{\textwidth}{!}{\documentclass[11pt, tikz, border=10pt]{standalone}
\usepackage{pgfplots}
\pgfplotsset{compat=1.18}
\usepackage{amsmath}
\usetikzlibrary{calc, arrows.meta, positioning, shapes.geometric, fit, patterns}

\definecolor{bglight}{HTML}{F8FAFC}
\definecolor{borderc}{HTML}{E2E8F0}
\definecolor{textmain}{HTML}{0F172A}
\definecolor{textsub}{HTML}{475569}
\definecolor{chartblue}{HTML}{0284C7}
\definecolor{queryred}{HTML}{EF4444}
\definecolor{valleygreen}{RGB}{35,95,25}
\definecolor{customorange}{RGB}{255,130,0}

\begin{document}

\begin{tikzpicture}[
    line join=round, 
    line cap=round,
    declare function={
        % Linear interpolation step (100% numerically stable, zero risk of TeX overflow)
        sstep(\x,\a,\b) = (\x < \a ? 0 : (\x > \b ? 1 : (\x - \a)/(\b - \a)));
        % Exact mathematical formula from page2.tex representing the user's skyline
        user_skyline(\x) = 0.35 + 0.20*sin(\x*200 + 45) + 0.08*cos(\x*380);
        % Base Curve 1 components (matching the 1.6 scale centered at 8.6 with new vertical offset)
        target_match(\x) = 5.90 + 1.6 * user_skyline((\x - 8.6) * 0.75);
        target_noise(\x) = 6.6 + 0.15*sin(\x*110) - 0.08*cos(\x*230);
        % Unified single continuous function with smooth transitions
        base_curve_one(\x) = (\x < 7.8 ? target_noise(\x) : (\x < 8.4 ? (1.0 - sstep(\x, 7.8, 8.4))*target_noise(\x) + sstep(\x, 7.8, 8.4)*target_match(\x) : (\x < 11.2 ? target_match(\x) : (\x < 11.7 ? (1.0 - sstep(\x, 11.2, 11.7))*target_match(\x) + sstep(\x, 11.2, 11.7)*target_noise(\x) : target_noise(\x) ))));
        % DTW Wave-form signals for the axes (derived from the real user's skyline)
        tsine(\x) = 1.2 * user_skyline(\x * 0.6) - 0.4;
        lsine(\x) = 1.2 * user_skyline(\x * 0.5 + 0.1) - 0.4;
    },
    headerstyle/.style={font=\large\sffamily\bfseries, text=textmain},
    substyle/.style={font=\small\sffamily\bfseries, text=textsub},
    boxstyle/.style={draw=borderc, thick, fill=white, rounded corners=4pt, inner sep=6pt, align=center},
    mathstyle/.style={draw=borderc, thick, fill=bglight, rounded corners=3pt, inner sep=6pt, align=center, font=\sffamily\small},
    arrowstyle/.style={-{Stealth[length=3mm, width=2mm]}, line width=1.5pt, draw=textsub}
]

    % ========================================================
    % 1. PHASE 1: VECTOR INDEX GRID & SLIDING WINDOW SEARCH
    % ========================================================
    \node[headerstyle] (title1) at (9.6, 13.2) {Phase 1: FFT-Accelerated Sliding Window ($L_2$) Search};
    
    % Vector Index Grid
    \begin{scope}[shift={(-5.0, 10.1)}]
        \draw[borderc, thick, fill=white, rounded corners=4pt] (0,0) rectangle (3.0,2.5);
        \node[font=\sffamily\bfseries\small, text=textmain] at (1.5, 2.1) {Vector Index};
        \draw[borderc, thin, step=0.3cm] (0.4,0.2) grid (2.5,1.7);
    \end{scope}

    % Skyline Extracted Box
    \begin{scope}[shift={(-5.0, 6.7)}]
        \draw[borderc, thick, fill=white, rounded corners=4pt] (0,0) rectangle (3.0,2.5);
        \node[font=\sffamily\bfseries\scriptsize, text=textmain, align=center] at (1.5, 2.1) {Skyline Extracted\\from User's Photo ($Q$)};
        % Exact page2.tex curve scaled to fit the display bounds beautifully
        \draw[queryred, line width=1.5pt] plot[domain=0.3:2.7, samples=40] (\x, {0.45 + 1.6*user_skyline((\x-0.3)*0.75)});
    \end{scope}

    % Mathematical Formula Box with clean, spacious derivation
    \node[mathstyle, text width=6.8cm, inner sep=10pt] (p1math) at (-3.5, 3.2) {
        \textbf{Distance Equation ($L_2$)}\\
        \vspace{4pt}
        $\begin{aligned}
            D_i^2 &= \sum_{k=1}^{m} (t_{i+k-1} - q_k)^2 \\
                  &= \sum_{k=1}^{m} q_k^2 + \sum_{k=1}^{m} t_{i+k-1}^2 - 2 d_i
        \end{aligned}$\\
        \vspace{8pt}
        \textbf{Cross-Correlation via FFT}\\
        \vspace{4pt}
        $\begin{aligned}
            d_i &= \sum_{k=1}^{m} t_{i+k-1} q_k \\
            d_i &= \left[\mathcal{F}^{-1}\left(\mathcal{F}(T) \odot \mathcal{F}(\tilde{Q})\right)\right]_{i+m-1}
        \end{aligned}$\\
        \vspace{8pt}
        \scriptsize\color{textsub} Where: $T$ is Database (len $N$), $Q$ is Query (len $m$).
    };

    % Sliding Window Curves Panel with expanded vertical padding (Height = 7.6)
    \begin{scope}[shift={(3.5, 5.2)}]
        % Background panel
        \draw[borderc, thick, fill=bglight, rounded corners=6pt] (0.0, 0.0) rectangle (12.2, 7.6);
        
        % ----------------------------------------------------
        % ROW 1: BASE CURVE 1 (Single continuous C1-smooth curve)
        % ----------------------------------------------------
        \draw[textsub!60, line width=1.5pt] plot[domain=0.5:11.7, samples=100] (\x, {base_curve_one(\x)});
        
        % Discrete dots perfectly aligned along the continuous curve
        \foreach \xd in {0.8,1.4,...,11.6} { 
            \fill[textmain] (\xd, {base_curve_one(\xd)}) circle (1.3pt); 
        }

        % ----------------------------------------------------
        % ROW 2: SLIDING WINDOWS (Under Base Curve 1)
        % ----------------------------------------------------
        % Window W_11 (Bad Match - Left)
        \draw[queryred, thick, rounded corners=4pt, fill=white, fill opacity=0.85] (1.0, 3.9) rectangle (3.8, 5.5);
        \node[font=\sffamily\tiny\bfseries, text=queryred, anchor=south] at (2.4, 5.55) {\textbf{Bad Match}};
        \draw[queryred, line width=1.5pt] plot[domain=1.2:3.6, samples=40] (\x, {3.95 + 1.6 * user_skyline((\x - 1.2) * 0.75)});
        
        % Projection lines from Row 1 down to Row 2 (W11)
        \foreach \xv in {1.6, 2.3, 3.0} {
            \draw[queryred, dashed, thin] (\xv, {base_curve_one(\xv)}) -- (\xv, {3.95 + 1.6 * user_skyline((\xv - 1.2) * 0.75)});
            \fill[queryred] (\xv, {3.95 + 1.6 * user_skyline((\xv - 1.2) * 0.75)}) circle (1.2pt);
        }

        \node[font=\large] at (5.9, 4.3) {\dots};

        % Window W_12 (Good Match - Right, sharing the exact terrain function)
        \draw[valleygreen, thick, rounded corners=4pt, fill=white, fill opacity=0.85] (8.4, 3.9) rectangle (11.2, 5.5);
        \node[font=\sffamily\tiny\bfseries, text=valleygreen, anchor=south] at (9.8, 5.55) {\textbf{Good Match}};
        \draw[queryred, line width=1.5pt] plot[domain=8.6:11.0, samples=40] (\x, {3.95 + 1.6 * user_skyline((\x - 8.6) * 0.75)});
        
        % Projection lines from Row 1 down to Row 2 (W12)
        \foreach \xv in {9.0, 9.7, 10.4} {
            \draw[valleygreen, dashed, thin] (\xv, {base_curve_one(\xv)}) -- (\xv, {3.95 + 1.6 * user_skyline((\xv - 8.6) * 0.75)});
            \fill[valleygreen] (\xv, {3.95 + 1.6 * user_skyline((\xv - 8.6) * 0.75)}) circle (1.2pt);
        }

        % ----------------------------------------------------
        % ROW 3: BASE CURVE 2 (Continuous, nicely separated)
        % ----------------------------------------------------
        \draw[textsub!60, line width=1.5pt] plot[domain=0.5:11.7, samples=60] (\x, {2.9 + 0.22*sin(\x*130) - 0.12*cos(\x*180) + 0.08*sin(\x*310)});
        \foreach \xd in {0.8,1.4,...,11.6} { \fill[textmain] (\xd, {2.9 + 0.22*sin(\xd*130) - 0.12*cos(\xd*180) + 0.08*sin(\xd*310)}) circle (1.3pt); }

        % ----------------------------------------------------
        % ROW 4: SLIDING WINDOWS (Under Base Curve 2)
        % ----------------------------------------------------
        % Window W_21 (Bad Match - Left)
        \draw[queryred, thick, rounded corners=4pt, fill=white, fill opacity=0.85] (1.0, 0.2) rectangle (3.8, 1.8);
        \node[font=\sffamily\tiny\bfseries, text=queryred, anchor=south] at (2.4, 1.85) {\textbf{Bad Match}};
        \draw[queryred, line width=1.5pt] plot[domain=1.2:3.6, samples=40] (\x, {0.25 + 1.6 * user_skyline((\x - 1.0) * 0.75)});
        
        % Projection lines from Row 3 down to Row 4 (W21)
        \foreach \xv in {1.6, 2.3, 3.0} {
            \draw[queryred, dashed, thin] (\xv, {2.9 + 0.22*sin(\xv*130) - 0.12*cos(\xv*180) + 0.08*sin(\xv*310)}) -- (\xv, {0.25 + 1.6 * user_skyline((\xv - 1.0) * 0.75)});
            \fill[queryred] (\xv, {0.25 + 1.6 * user_skyline((\xv - 1.0) * 0.75)}) circle (1.2pt);
        }

        \node[font=\large] at (5.9, 0.8) {\dots};

        % Window W_22 (Bad Match - Right)
        \draw[queryred, thick, rounded corners=4pt, fill=white, fill opacity=0.85] (8.4, 0.2) rectangle (11.2, 1.8);
        \node[font=\sffamily\tiny\bfseries, text=queryred, anchor=south] at (9.8, 1.85) {\textbf{Bad Match}};
        \draw[queryred, line width=1.5pt] plot[domain=8.6:11.0, samples=40] (\x, {0.25 + 1.6 * user_skyline((\x - 8.6) * 0.75)});
        
        % Projection lines from Row 3 down to Row 4 (W22)
        \foreach \xv in {9.0, 9.7, 10.4} {
            \draw[queryred, dashed, thin] (\xv, {2.9 + 0.22*sin(\xv*130) - 0.12*cos(\xv*180) + 0.08*sin(\xv*310)}) -- (\xv, {0.25 + 1.6 * user_skyline((\xv - 8.6) * 0.75)});
            \fill[queryred] (\xv, {0.25 + 1.6 * user_skyline((\xv - 8.6) * 0.75)}) circle (1.2pt);
        }
    \end{scope}

    % Two consistent dark-gray arrows pointing from Vector Index to the base curves
    \draw[arrowstyle] (-2.0, 11.35) -- (3.3, 11.8); % Arrow to Top Base Curve
    \draw[arrowstyle] (-2.0, 10.9) to[out=-35, in=180] (3.3, 8.1); % Curved Arrow to Bottom Base Curve

    % Vertical ellipsis positioned perfectly in the whitespace between the two arrows
    \node[font=\Large, text=textsub!60] at (0.6, 10.5) {$\vdots$};

    % Flowchart arrow connecting Phase 1 and Phase 2 (Clean vertical connection)
    \draw[arrowstyle] (9.6, 4.5) -- (9.6, 1.2);

    % Phase 1 Candidate Locations Table
    \begin{scope}[shift={(16.2, 7.1)}]
        \node[font=\sffamily\bfseries\small, text=textmain, anchor=south] at (3.25, 3.8) {Probable Locations};
        \draw[borderc, thick, fill=white] (0,0) rectangle (6.5, 3.8);
        \draw[borderc, thin] (0, 0.76) -- (6.5, 0.76);
        \draw[borderc, thin] (0, 1.52) -- (6.5, 1.52);
        \draw[borderc, thin] (0, 2.28) -- (6.5, 2.28);
        \draw[borderc, thin] (0, 3.04) -- (6.5, 3.04);
        \draw[borderc, thin] (1.2, 0) -- (1.2, 3.8);
        \draw[borderc, thin] (2.8, 0) -- (2.8, 3.8);
        \draw[borderc, thin] (4.4, 0) -- (4.4, 3.8);
        
        \node[font=\sffamily\bfseries\scriptsize] at (0.6, 3.42) {Rank};
        \node[font=\sffamily\bfseries\scriptsize] at (2.0, 3.42) {Latitude};
        \node[font=\sffamily\bfseries\scriptsize] at (3.6, 3.42) {Longitude};
        \node[font=\sffamily\bfseries\scriptsize] at (5.45, 3.42) {Distance};
        
        \node[font=\sffamily\scriptsize] at (0.6, 2.66) {1};
        \node[font=\sffamily\scriptsize] at (2.0, 2.66) {$Lat_1$};
        \node[font=\sffamily\scriptsize] at (3.6, 2.66) {$Lon_1$};
        \node[font=\sffamily\scriptsize\bfseries, text=valleygreen] at (5.45, 2.66) {$L_2$};

        \node[font=\sffamily\scriptsize] at (0.6, 1.9) {2};
        \node[font=\sffamily\scriptsize] at (2.0, 1.9) {$Lat_2$};
        \node[font=\sffamily\scriptsize] at (3.6, 1.9) {$Lon_2$};
        \node[font=\sffamily\scriptsize\bfseries, text=queryred] at (5.45, 1.9) {$L_2$};

        \node[font=\sffamily\scriptsize] at (0.6, 1.14) {$\vdots$};
        \node[font=\sffamily\scriptsize] at (2.0, 1.14) {$\vdots$};
        \node[font=\sffamily\scriptsize] at (3.6, 1.14) {$\vdots$};
        \node[font=\sffamily\scriptsize] at (5.45, 1.14) {$\vdots$};

        \node[font=\sffamily\scriptsize] at (0.6, 0.38) {100};
        \node[font=\sffamily\scriptsize] at (2.0, 0.38) {$Lat_N$};
        \node[font=\sffamily\scriptsize] at (3.6, 0.38) {$Lon_N$};
        \node[font=\sffamily\scriptsize\bfseries, text=queryred] at (5.45, 0.38) {$L_2$};
    \end{scope}

    % ========================================================
    % 2. PHASE 2: SAKOE-CHIBA CONSTRAINED DTW ALIGNMENT
    % ========================================================
    \node[headerstyle] (title2) at (9.6, 0.6) {Phase 2: Sakoe-Chiba Constrained DTW Alignment};

    % DTW Recurrence Formula (Shifted down to -3.8 to eliminate left column overlap)
    \node[mathstyle, text width=5.5cm] (p2math) at (-3.5, -3.8) {
        \textbf{DTW Recurrence Formula}\\
        \vspace{4pt}
        $D(i, j) = C(i, j) + \min \begin{cases} D(i-1, j) \\ D(i, j-1) \\ D(i-1, j-1) \end{cases}$
    };

    \begin{scope}[shift={(3.5, -4.5)}]
% Background panel with expanded padding at top (y = 4.6) and bottom (y = -2.8)
        \draw[borderc, thick, fill=bglight, rounded corners=6pt] (0.0, -2.8) rectangle (12.2, 4.6);

        % DTW Matrix 1: Full warping path inside corridor (Top-Left to Bottom-Right)
        \begin{scope}[shift={(1.2, 0.6)}]
            % Shaded Sakoe-Chiba Corridor sloping from Top-Left to Bottom-Right
            \foreach \x/\y in {0/9, 0/8, 1/9, 1/8, 1/7, 2/8, 2/7, 2/6, 3/7, 3/6, 3/5, 4/6, 4/5, 4/4, 5/5, 5/4, 5/3, 6/4, 6/3, 6/2, 7/3, 7/2, 7/1, 8/2, 8/1, 8/0, 9/1, 9/0} {
                \fill[chartblue!15] (\x*0.3, \y*0.3) rectangle (\x*0.3+0.3, \y*0.3+0.3);
            }
            \draw[borderc, thin, step=0.3cm] (0,0) grid (3.0,3.0);
            
            % Correctly oriented and positioned Band label
            \node[font=\sffamily\tiny\bfseries, text=chartblue, rotate=-45] at (1.5, 1.7) {Sakoe-Chiba Band};

            % Completed Path (represented by gray dots & connections)
            \draw[lightgray, thick] (0.15, 2.85) -- (0.45, 2.55) -- (0.45, 2.25) -- (0.75, 1.95) -- (1.05, 1.65) -- (1.35, 1.35) -- (1.65, 1.05) -- (1.95, 0.75) -- (2.25, 0.75) -- (2.55, 0.45) -- (2.85, 0.15);
            \foreach \cx/\cy in {0.15/2.85, 0.45/2.55, 0.45/2.25, 0.75/1.95, 1.05/1.65, 1.35/1.35, 1.65/1.05, 1.95/0.75, 2.25/0.75, 2.55/0.45, 2.85/0.15} {
                \fill[lightgray] (\cx, \cy) circle (2.2pt);
            }

            % Highlighted warping path node
            \fill[black] (1.05, 1.65) circle (2.8pt);

            % Wave-form signals plotted along the horizontal axis
            \draw[thick, ->] (-0.1, 3.3) -- (3.2, 3.3);
            \draw[customorange, line width=1.5pt] plot[domain=0:3.0, samples=40] (\x, {3.3 + tsine(\x)});
            \foreach \ix in {0,1,...,9} {
                \pgfmathsetmacro{\currx}{\ix*0.3 + 0.15}
                \fill[customorange] (\currx, {3.3 + tsine(\currx)}) circle (1.8pt);
            }

            % Wave-form signals plotted along vertical axis
            \draw[thick, ->] (-0.3, 3.1) -- (-0.3, -0.2);
            \draw[customorange, line width=1.5pt] plot[domain=0:3.0, samples=40] ({-0.3 - lsine(\x)}, \x);
            \foreach \iy in {0,1,...,9} {
                \pgfmathsetmacro{\curry}{\iy*0.3 + 0.15}
                \fill[customorange] ({-0.3 - lsine(\curry)}, \curry) circle (1.8pt);
            }

            % Alignment projections intersecting at the key matrix node
            \draw[dashed, thin, black] (1.05, {3.3 + tsine(1.05)}) -- (1.05, 1.65);
            \draw[dashed, thin, black] ({-0.3 - lsine(1.65)}, 1.65) -- (1.05, 1.65);
        \end{scope}

        \node[font=\large] at (5.9, 2.1) {\dots};

        % DTW Matrix 2: Early Abandonment & Pruning (Label cleanly integrated at the top)
        \begin{scope}[shift={(7.6, 0.6)}]
            % Shaded Sakoe-Chiba Corridor sloping from Top-Left to Bottom-Right
            \foreach \x/\y in {0/9, 0/8, 1/9, 1/8, 1/7, 2/8, 2/7, 2/6, 3/7, 3/6, 3/5, 4/6, 4/5, 4/4, 5/5, 5/4, 5/3, 6/4, 6/3, 6/2, 7/3, 7/2, 7/1, 8/2, 8/1, 8/0, 9/1, 9/0} {
                \fill[chartblue!15] (\x*0.3, \y*0.3) rectangle (\x*0.3+0.3, \y*0.3+0.3);
            }
            \draw[borderc, thin, step=0.3cm] (0,0) grid (3.0,3.0);

            % Early Terminated path
            \draw[lightgray, thick] (0.15, 2.85) -- (0.45, 2.55) -- (0.45, 2.25) -- (0.75, 1.95);
            \foreach \cx/\cy in {0.15/2.85, 0.45/2.55, 0.45/2.25, 0.75/1.95} {
                \fill[lightgray] (\cx, \cy) circle (2.2pt);
            }
            
            % Abandonment cross positioned at (1.05, 1.65)
            \node[queryred, scale=1.3] at (1.05, 1.65) {$\mathbf{\times}$};

            % Axes signals
            \draw[thick, ->] (-0.1, 3.3) -- (3.2, 3.3);
            \draw[customorange, line width=1.5pt] plot[domain=0:3.0, samples=40] (\x, {3.3 + tsine(\x)});
            \foreach \ix in {0,1,...,9} {
                \pgfmathsetmacro{\currx}{\ix*0.3 + 0.15}
                \fill[customorange] (\currx, {3.3 + tsine(\currx)}) circle (1.8pt);
            }

            \draw[thick, ->] (-0.3, 3.1) -- (-0.3, -0.2);
            \draw[customorange, line width=1.5pt] plot[domain=0:3.0, samples=40] ({-0.3 - lsine(\x)}, \x);
            \foreach \iy in {0,1,...,9} {
                \pgfmathsetmacro{\curry}{\iy*0.3 + 0.15}
                \fill[customorange] ({-0.3 - lsine(\curry)}, \curry) circle (1.8pt);
            }

            % Cleanly integrated label safely at the top inside the expanded panel
            \node[font=\sffamily\bfseries\scriptsize, text=queryred] at (1.5, 3.89) {Early Abandoning \& Pruning};
        \end{scope}

        % Elastic Horizon Alignment Visualizer (Drawn using the real user double-bump skyline)
        \begin{scope}[shift={(1.0, -2.4)}]
            \draw[borderc, thick, fill=white] (0.2, 0.4) rectangle (9.8, 2.8);

            % Standardized orange curves mapping elastic points
            \draw[customorange, line width=1.5pt] plot[domain=0.5:9.5, samples=50] (\x, {2.0 + 0.8 * user_skyline((\x - 0.5) * 0.2)});
            \draw[customorange, line width=1.5pt] plot[domain=0.5:9.5, samples=50] (\x, {0.8 + 0.8 * user_skyline((\x - 0.5) * 0.17 + 0.1)});

            \foreach \ix in {0,1,...,18} {
                \pgfmathsetmacro{\cx}{0.5 + \ix*0.5}
                \pgfmathsetmacro{\targetx}{0.5 + \ix*0.5 - 0.35*sin(\ix*10)}
                \draw[textsub!45, thin] (\cx, {2.0 + 0.8 * user_skyline((\cx - 0.5) * 0.2)}) -- (\targetx, {0.8 + 0.8 * user_skyline((\targetx - 0.5) * 0.17 + 0.1)});
                \fill[customorange] (\cx, {2.0 + 0.8 * user_skyline((\cx - 0.5) * 0.2)}) circle (1.8pt);
                \fill[customorange] (\targetx, {0.8 + 0.8 * user_skyline((\targetx - 0.5) * 0.17 + 0.1)}) circle (1.8pt);
            }

            % Highlighted optimal warping path connection
            \pgfmathsetmacro{\cx}{0.5 + 9*0.5}
            \pgfmathsetmacro{\targetx}{0.5 + 9*0.5 - 0.35*sin(90)}
            \draw[black, very thick] (\cx, {2.0 + 0.8 * user_skyline((\cx - 0.5) * 0.2)}) -- (\targetx, {0.8 + 0.8 * user_skyline((\targetx - 0.5) * 0.17 + 0.1)});
        \end{scope}
    \end{scope}

    % Phase 2 Candidate Locations Table
    \begin{scope}[shift={(16.2, -5.45)}]
        \node[font=\sffamily\bfseries\small, text=textmain, anchor=south] at (3.25, 3.8) {Probable Locations};
        \draw[borderc, thick, fill=white] (0,0) rectangle (6.5, 3.8);
        \draw[borderc, thin] (0, 0.76) -- (6.5, 0.76);
        \draw[borderc, thin] (0, 1.52) -- (6.5, 1.52);
        \draw[borderc, thin] (0, 2.28) -- (6.5, 2.28);
        \draw[borderc, thin] (0, 3.04) -- (6.5, 3.04);
        \draw[borderc, thin] (1.2, 0) -- (1.2, 3.8);
        \draw[borderc, thin] (2.8, 0) -- (2.8, 3.8);
        \draw[borderc, thin] (4.4, 0) -- (4.4, 3.8);
        
        \node[font=\sffamily\bfseries\scriptsize] at (0.6, 3.42) {Rank};
        \node[font=\sffamily\bfseries\scriptsize] at (2.0, 3.42) {Latitude};
        \node[font=\sffamily\bfseries\scriptsize] at (3.6, 3.42) {Longitude};
        \node[font=\sffamily\bfseries\scriptsize] at (5.45, 3.42) {Cost};
        
        \node[font=\sffamily\scriptsize] at (0.6, 2.66) {1};
        \node[font=\sffamily\scriptsize] at (2.0, 2.66) {$Lat_1$};
        \node[font=\sffamily\scriptsize] at (3.6, 2.66) {$Lon_1$};
        \node[font=\sffamily\scriptsize\bfseries, text=valleygreen] at (5.45, 2.66) {DTW};

        \node[font=\sffamily\scriptsize] at (0.6, 1.9) {2};
        \node[font=\sffamily\scriptsize] at (2.0, 1.9) {$Lat_2$};
        \node[font=\sffamily\scriptsize] at (3.6, 1.9) {$Lon_2$};
        \node[font=\sffamily\scriptsize\bfseries, text=valleygreen] at (5.45, 1.9) {DTW};

        \node[font=\sffamily\scriptsize] at (0.6, 1.14) {$\vdots$};
        \node[font=\sffamily\scriptsize] at (2.0, 1.14) {$\vdots$};
        \node[font=\sffamily\scriptsize] at (3.6, 1.14) {$\vdots$};
        \node[font=\sffamily\scriptsize\bfseries, text=valleygreen] at (5.45, 1.14) {DTW};

        \node[font=\sffamily\scriptsize] at (0.6, 0.38) {5};
        \node[font=\sffamily\scriptsize] at (2.0, 0.38) {$Lat_5$};
        \node[font=\sffamily\scriptsize] at (3.6, 0.38) {$Lon_5$};
        \node[font=\sffamily\scriptsize\bfseries, text=valleygreen] at (5.45, 0.38) {DTW};
    \end{scope}

\end{tikzpicture}
\end{document}
}
    \caption{Detailed architecture of the matching engine.}
    \label{fig:matching_engine}
\end{figure}

\section{Validation}
We evaluate the system in two ways. On the synthetic side, we hold out a subset of generated horizons and measure localization accuracy on this test split. For real-world evaluation, we use geotagged panoramas from Google Street View's Everest trail coverage \cite{google2015streetview}. The Khumbu trails have sufficient Street View data for this purpose. Rather than matching directly against full panoramas, we decompose each one into standard perspective photographs and run the pipeline on these individual frames, comparing the predicted coordinates against the known GPS tags.
